\chapter{CƠ SỞ LÝ THUYẾT}

\section{Công nghệ AI và DL}

\subsection{Giới thiệu về DL}

DL là một nhánh quan trọng của Machine Learning, sử dụng các mạng nơ-ron nhân tạo nhiều lớp để tự động học và trích xuất đặc trưng từ dữ liệu. 
Trong BarberGo, DL được áp dụng cho:
\begin{itemize}
    \item CNN: Phân loại mụn (MobileNetV2)
    \item Diffusion Models: Tạo kiểu tóc (SD1.5)
\end{itemize}

\subsection{Mô hình sinh ảnh Generative AI – SD1.5}

SD1.5 là một mô hình Generative AI dựa trên Diffusion Model, có khả năng sinh ảnh từ mô tả văn bản (text-to-image) và chỉnh sửa ảnh (image inpainting). Mô hình hoạt động bằng cách dần loại bỏ nhiễu khỏi dữ liệu ngẫu nhiên để tạo ra hình ảnh phù hợp với prompt đầu vào.

Kiến trúc của SD1.5 bao gồm các thành phần chính:
\begin{itemize}
    \item Text Encoder (CLIP): Chuyển đổi prompt văn bản thành vector đặc trưng ngữ nghĩa.
    \item U-Net: Thực hiện quá trình khử nhiễu trong không gian tiềm ẩn (latent space).
    \item Variational Autoencoder (VAE): Mã hóa ảnh vào latent space và giải mã thành ảnh output.
\end{itemize}

\subsubsection{SD1.5 Inpainting trong BarberGo}

Pipeline xử lý:
\begin{enumerate}
    \item \textbf{Phát hiện khuôn mặt:} MediaPipe tạo hair mask và face protection mask.
    \item \textbf{Inpainting:} U-Net thay đổi vùng tóc theo prompt, bảo vệ khuôn mặt. 
          Số bước: 30-40 steps, guidance scale: 7.5-8.0.
    \item \textbf{Output:} VAE Decoder sinh ảnh 512×512, thời gian ~10-15 giây/ảnh.
\end{enumerate}

Ưu điểm: Chỉ thay đổi tóc, giữ nguyên mặt; linh hoạt với nhiều prompt; kết quả tự nhiên.

\subsection{MediaPipe trong nhận diện và phân vùng khuôn mặt}

MediaPipe là thư viện mã nguồn mở do Google phát triển, hỗ trợ xây dựng các pipeline xử lý hình ảnh và video theo thời gian thực. MediaPipe Face Mesh cung cấp 468 điểm mốc (landmarks) trên khuôn mặt với độ chính xác cao, cho phép xác định chính xác các vùng như trán, má, cằm, mắt, mũi và miệng.

\subsection{MediaPipe}

MediaPipe Face Mesh cung cấp 468 landmarks để:
\begin{itemize}
    \item Phân vùng da: 5 vùng (trán, má trái/phải, mũi, cằm) cho Acne Detection
    \item Tạo hair mask cho Hairstyle Generation
\end{itemize}

Ưu điểm: Nhanh (<100ms), chạy trên CPU, độ chính xác cao.

\subsection{CNN}

CNN là kiến trúc mạng nơ-ron chuyên xử lý dữ liệu dạng lưới như hình ảnh. CNN có khả năng tự động trích xuất đặc trưng quan trọng thông qua các lớp tích chập (convolution) và gộp (pooling), giúp giảm số lượng tham số và tăng hiệu quả học.

\subsubsection{Kiến trúc cơ bản}

1. Lớp tích chập (Convolutional Layer):
\begin{equation}
Y_{i,j} = \sigma\left(\sum_{m}\sum_{n} W_{m,n} \cdot X_{i+m,j+n} + b\right)
\end{equation}
Trong đó: $W$ là kernel (bộ lọc), $b$ là bias, $\sigma$ là activation function.

2. Lớp gộp (Pooling Layer):
\begin{equation}
Y_{i,j} = \max_{m,n \in R} X_{i+m,j+n}
\end{equation}
Max pooling lấy giá trị lớn nhất trong vùng $R$ để giảm kích thước feature map.

3. Lớp kết nối đầy đủ (Fully Connected):
\begin{equation}
y = \sigma(W \cdot x + b)
\end{equation}

4. Hàm Softmax (phân loại):
\begin{equation}
P(y=k) = \frac{e^{z_k}}{\sum_{j=1}^{K} e^{z_j}}
\end{equation}

5. Hàm mất mát Binary Cross-Entropy:
\begin{equation}
\mathcal{L} = -\frac{1}{N} \sum_{i=1}^{N} \left[y_i \log(\hat{y}_i) + (1-y_i) \log(1-\hat{y}_i)\right]
\end{equation}

\subsubsection{CNN cho Acne Detection}

Model: MobileNetV2 (pretrained) với binary classification (có mụn / không mụn)
\begin{itemize}
    \item Input: 224×224 RGB (5 vùng da từ MediaPipe)
    \item Preprocessing: Resize, normalize, Gaussian Blur (bỏ lỗ chân lông)
    \item Training: Binary Cross-Entropy, Adam optimizer, 50 epochs
    \item Output: Confidence score [0,1], phân loại severity (mild/moderate/severe)
\end{itemize}

\subsection{RAG và Chatbot}

RAG là kỹ thuật kết hợp giữa hệ thống truy xuất thông tin (Information Retrieval) và mô hình NLG. 
RAG cho phép chatbot truy xuất tri thức từ cơ sở dữ liệu trước khi sinh câu trả lời, giúp tăng độ chính xác, giảm hallucination (bịa đặt thông tin) và kiến thức lõi thời.

\subsubsection{Kiến trúc RAG trong BarberGo}

Hệ thống RAG chatbot bao gồm 3 thành phần chính:

\begin{enumerate}
    \item \textbf{Embedding Model}: Google text-embedding-004, chuyển văn bản thành vector 768 chiều với mục đích trích xuất đặc trưng ngữ nghĩa.

    \item \textbf{Vector Database}: Supabase PostgreSQL + pgvector extension, lưu trữ embeddings và tìm kiếm semantic với cosine similarity với mục đích truy xuất tri thức mà BarberGo có.

    \item \textbf{Generative Model}: Google Gemini 2.0 Flash là Multimodel LLM cung cấp qua Generative AI API, sinh câu trả lời dựa trên context trích xuất từ vector database.
\end{enumerate}

\textbf{Quy trình chi tiết gồm 3 phần chính:}
\begin{enumerate}
    \item Người dùng nhập câu hỏi.
    \item Hệ thống tạo embedding vector (768-dim) cho câu hỏi.
    \item Tìm kiếm top-K documents tương tự nhất trong vector database (cosine similarity > 0.5).
    \item Trích xuất context từ documents và kết hợp với câu hỏi.
    \item Gemini API sinh câu trả lời dựa trên context.
    \item Đánh giá confidence (high/medium/low) dựa trên similarity score.
    \item Trả kết quả cho người dùng và lưu vào chat history.
\end{enumerate}
\textbf{Ứng dụng trong BarberGo:}
\begin{itemize}
    \item Tư vấn dịch vụ cắt tóc, chăm sóc da.
    \item Hướng dẫn sử dụng app (đặt lịch, hủy lịch, hợp tác).
    \item Giải đáp thắc mắc về salon, giá cả, thời gian.
    \item Hỗ trợ khách hàng 24/7 tự động.
\end{itemize}


\section{Công nghệ Backend}

\subsection{Python}

Backend của BarberGo được xây dựng bằng python với FastAPI. Python được sử dụng bởi vì:
\begin{itemize}
    \item Hỗ trợ AI/ML libraries (Diffusers, PyTorch, MediaPipe)
    \item Tích hợp dễ Google Gemini API và Supabase SDK
    \item Cộng đồng lớn, tài liệu phong phú
\end{itemize}

\subsection{FastAPI}

FastAPI là ASGI framework hiệu năng cao, hỗ trợ async/await natively. 
Ưu điểm:
\begin{itemize}
    \item Xử lý concurrent AI requests hiệu quả
    \item Auto validation với Pydantic
    \item Auto generate API docs (Swagger UI)
\end{itemize}

\section{Công nghệ Frontend}
\subsection{Flutter}

Flutter (Dart, Google) cho phép xây dựng mobile app (iOS/Android) từ single codebase. 
Ưu điểm: Hot reload nhanh, UI đẹp (60fps), package ecosystem phong phú.

\subsection{Vue.js 3}

Vue.js 3 với Composition API cho web admin dashboard. 
Stack: Vite (build tool), Pinia (state management), TailwindCSS (styling).

\subsection{Figma}
Figma là công cụ thiết kế UI/UX cloud-based, hỗ trợ collaboration thời gian thực. Trong BarberGo, Figma được sử dụng để thiết kế wireframes, mockups cho mobile app và admin dashboard, xây dựng design system (colors, typography, components), tạo interactive prototypes và export assets cho developers.


\section{Cơ sở dữ liệu}
\subsection{Supabase}

Supabase (BaaS) cung cấp:
\begin{itemize}
    \item PostgreSQL database với pgvector, PostGIS extensions
    \item Authentication (Email, OAuth)
    \item Storage (S3-compatible)
    \item Row Level Security (RLS)
\end{itemize}

Ưu điểm: Setup nhanh, free tier generous, auto-generated REST APIs.

\subsection{Pgvector Extension}

Pgvector là PostgreSQL extension cho vector similarity search, được sử dụng chủ yếu cho Retrieval-Augmented Generation (RAG) và semantic search.


\section{Cloud Services}

\subsection{Google Cloud Platform}

Google Cloud Platform (GCP) là dịch vụ cloud computing của Google, cung cấp infrastructure và AI/ML services.

\textbf{Services sử dụng trong BarberGo:}

\begin{enumerate}
    \item \textbf{Gemini AI API:}
    \begin{itemize}
        \item Model: Gemini 2.0 Flash (fast, cost-effective).
        \item Text generation cho RAG chatbot.
        \item Embedding generation: text-embedding-004 (768-dim).
        \item API endpoint: \texttt{generativelanguage.googleapis.com}.
    \end{itemize}

    \item \textbf{Firebase Cloud Messaging (FCM):}
    \begin{itemize}
        \item Push notifications cho mobile app.
        \item Cross-platform (iOS, Android).
        \item Topics, device groups.
        \item Delivery reports, analytics.
    \end{itemize}
\end{enumerate}

\subsection{Supabase Storage}

Supabase Storage được tích hợp sẵn trong Supabase.

\textbf{Chức năng:}
\begin{itemize}
    \item Upload/download files qua API.
    \item CDN distribution cho fast delivery.
    \item Public và private buckets.
    \item Row Level Security (RLS) cho access control.
    \item Image transformations (resize, crop).
\end{itemize}

\section{Kiến trúc hệ thống}

\textbf{Frontend (MVVM):} 
- Model: Data từ API
- View: UI components
- ViewModel: Business logic, state management

\textbf{Backend (Three-Tier):}
- Presentation: FastAPI routers, validation
- Business Logic: Services (RAG, Booking, Auth)
- Data Access: Supabase client, CRUD operations

\subsection{Geocoding và Maps}

Hệ thống sử dụng:
\begin{itemize}
    \item Nominatim API (OpenStreetMap) cho geocoding
    \item PostGIS để lưu trữ và truy vấn dữ liệu địa lý
    \item Routing algorithms: Dựa vào khoảng cách tiến hành lựa chọn thuật toán phù hợp: Dijkstra (<5km), A* (5-15km), Bellman-Ford (>15km)
\end{itemize}

\subsection{Postman Testing}
Postman được sử dụng để kiểm thử API backend:
\begin{itemize}
    \item Tạo Postman collection cho các endpoints
    \item Viết test scripts để validate responses
    \item Báo cáo lỗi và hiệu năng qua Postman Dashboard
\end{itemize}
\clearpage
